\section{Introdução}

A doença de Parkinson (\textit{Parkinson's disease} - PD) é uma enfermidade
neurodegenerativa progressiva cujo diagnóstico ainda depende de exame clínico e
escalas subjetivas, que podem se confundir com outros distúrbios do movimento
(\textit{differential diagnoses} - DD) \cite{varghese_machine_2024}. Por isso,
\textit{smartwatches} equipados com unidades de medição inercial
(\textit{inertial measurement unit} - IMU) vêm ganhando espaço na busca por
biomarcadores digitais mais objetivos \cite{ouabd_overview_2026}.

O lançamento do conjunto público \textit{Parkinson's Disease Smartwatch} (PADS)
\cite{varghese_machine_2024} reuniu em uma mesma base sinais inerciais dos dois
punhos, o questionário de sintomas não motores da doença de Parkinson (PDNMS) e
dados demográficos para PD, controles saudáveis (\textit{healthy controls} - HC)
e diagnósticos diferenciais (\textit{differential diagnoses} - DD). No entanto,
os trabalhos publicados desde então em aprendizado de máquina (\textit{machine
learning} - ML) e aprendizado profundo (\textit{deep learning} - DL) adotam
recortes amostrais, tarefas, particionamentos e métricas bastante distintos,
além de correrem risco de vazamento de dados (\textit{data leakage}) quando
sinais do mesmo participante ou etapas de balanceamento cruzam treino e teste
\cite{djemaa_quantifying_2026,pinto_neto_machine_2026}.

Neste artigo, conduzimos um mapeamento sistemático rápido
(\textit{rapid review}) e descritivo da literatura que utiliza o PADS. A partir
da estrutura \textit{Goal-Question-Metric} (GQM)
\pend{referência de GQM no Zotero}, buscamos caracterizar tarefas, recortes,
técnicas, validação, qualidade (QA), risco de vazamento relatado e lacunas
(QP1--QP5), reunindo \nTotal\ estudos (\nRevisadosPares\ revisados por pares e
\nPreprints\ \textit{preprints}) com análise de sensibilidade. O restante do
texto divide-se em Contexto (Seção~\ref{sec:contexto}), Método
(Seção~\ref{sec:metodo}), Resultados (Seção~\ref{sec:resultados}), Discussão
(Seção~\ref{sec:discussao}) e Conclusão (Seção~\ref{sec:conclusao}).
